\section{Docker Deployment}

\subsection{Architecture}

The Docker Compose setup runs two containers that share a tmux socket volume, allowing
the dashboard to read agent panes with zero modifications to the agent logic.

\begin{figure}[H]
\centering
\begin{tikzpicture}[
  box/.style={draw, rounded corners=4pt,
              fill=codebg, draw=codeframe, line width=1pt,
              text width=4.8cm, align=center, inner sep=10pt},
  vol/.style={draw=gray!70, fill=gray!15, rounded corners=3pt,
              text width=6.4cm, align=center, inner sep=6pt},
  arr/.style={-{Stealth[length=5pt]}, thick, gray!60}
]

% --- mas container ---
\node[box] (mas) at (0,0) {%
  {\small\bfseries\ttfamily mas}~{\footnotesize(ubuntu:22.04)}\\[5pt]
  {\footnotesize\ttfamily startmas.sh} {\footnotesize$\to$ tmux \texttt{f1\_race}}\\[3pt]
  {\footnotesize SICStus mounted from host (ro)}%
};

% --- ui container ---
\node[box, right=1.8cm of mas] (ui) {%
  {\small\bfseries\ttfamily ui}~{\footnotesize(python:3.11-slim)}\\[5pt]
  {\footnotesize\ttfamily dashboard.py}~{\footnotesize(Flask \texttt{:5000})}\\[3pt]
  {\footnotesize reads tmux panes via volume}%
};

% --- shared volume ---
\node[vol] (vol) at ($(mas.south)!0.5!(ui.south) + (0,-1.6)$) {%
  {\footnotesize\ttfamily tmux\_sock}~{\footnotesize volume\quad\texttt{/tmp/tmux-shared}}%
};

% --- arrows ---
\draw[arr] (mas.south) -- ++(0,-0.5) -| (vol.west);
\draw[arr] (ui.south)  -- ++(0,-0.5) -| (vol.east);

\end{tikzpicture}
\caption{Docker Compose architecture: the \texttt{mas} and \texttt{ui} containers share a tmux socket volume.}
\label{fig:docker_arch}
\end{figure}

\noindent
\textbf{Startup order} (managed by Docker healthchecks):
\begin{enumerate}[leftmargin=2em]
  \item \texttt{ui} starts immediately and is healthy once \texttt{/api/config} responds.
  \item \texttt{mas} waits for \texttt{ui} to be healthy, then runs \texttt{startmas.sh}.
\end{enumerate}

\subsection{Quickstart}
As in the native approach, the configuration requires:
\begin{enumerate}
  \item Clone the DALI repo from \url{https://github.com/AAAI-DISIM-UnivAQ/DALI} and navigate to the \texttt{DALI/Examples} folder.
  \item Clone the GitHub repo project from \url{https://github.com/Mik1810/f1_dali_project} under \texttt{Examples} folder.
\end{enumerate}

Then, execute the following commands:
\begin{lstlisting}[language=bash, numbers=none]
cd DALI/Examples/f1_race

# 1. Auto-detect SICStus and write .env  (run once)
bash docker/setup.sh

# 2. Build and start
docker compose up --build -d
\end{lstlisting}

The containers need a few seconds to start up and run the agents. Once ready, open \url{http://localhost:5000} in a browser to view the dashboard.

\subsection{Environment Variables}
The \texttt{mas} container needs the SICSTUS\_PATH environment variable
 to locate the SICStus installation. This is set in 
 \texttt{docker-compose.yml} and can be configured via a 
 \texttt{.env} file in the project root. 
This setup is done by the \texttt{docker/setup.sh} script, 
which auto-detects the SICStus path on the host and writes it 
to \texttt{.env}. In case of a different setup of the SICStus installation, 
the variable can be manually set in \texttt{.env} as follows:

\begin{longtable}{@{} l l p{7cm} @{}}
\toprule
\textbf{Variable} & \textbf{Default} & \textbf{Description} \\
\midrule
\endhead
\texttt{SICSTUS\_PATH} & \texttt{/usr/local/sicstus4.6.0} & Host path to SICStus install, mounted read-only into the \texttt{mas} container \\
\bottomrule
\end{longtable}
